\section{Endomorphism algebras and where Weil type sits}\label{app:albert}

\Cref{def:weiltype} asks for three things: an even dimension, an imaginary
quadratic field acting with balanced eigenvalue multiplicities, and
compatibility with the polarisation. This appendix places those three
conditions inside Albert's classification of endomorphism algebras. Two
statements come out of it that are used in the body. The first is that the
compatibility condition is automatic once the field is stable under the Rosati
involution, so Weil type really has only two conditions
(\Cref{lem:rosatiauto}). The second is that the balanced condition is not a
normalisation but the exact condition under which the Weil line carries a
Hodge class (\Cref{prop:signatureforced}), which is what makes the whole
subject about balanced signatures and nothing else.

\subsection{Positive involutions}%
\footnote{The numbering of the four types is Albert's and has stuck, although it corresponds to no natural order: types I to III have a totally real centre and type IV a CM one, so the natural division is two against two rather than three against one. Shimura's normalisation of the same classification, in terms of the associated symmetric domains, is the one that makes the moduli side transparent.}

\begin{definition}\label{def:positive}
Let $D$ be a finite dimensional semisimple $\QQ$-algebra. An involution
$\dagger$ of $D$, that is a $\QQ$-linear anti-automorphism with
$\dagger^{2}=\id$, is \emph{positive} if
\[
  \Tr_{D/\QQ}\bigl(x\,x^{\dagger}\bigr)>0
  \qquad\text{for every }x\in D,\ x\ne0,
\]
where $\Tr_{D/\QQ}$ is the trace of left multiplication on $D$ as a
$\QQ$-vector space.
\end{definition}

For an abelian variety $A$ with a polarisation $\eta$, the algebra
$D=\End^{0}(A)=\End(A)\otimes\QQ$ is semisimple and the Rosati involution
$\dagger$ attached to $\eta$ is positive \cite[Ch.~5]{BL04}. Every statement
below is a statement about the pair $(D,\dagger)$.

\begin{lemma}[Totally real centre]\label{lem:totallyreal}
Let $(D,\dagger)$ be as above with $D$ a division algebra, let $F$ be the
centre of $D$ and $F_{0}=\{z\in F:z^{\dagger}=z\}$. Then $F_{0}$ is a totally
real number field.
\end{lemma}

\begin{proof}
For $z\in F_{0}$ the positivity condition reads
$\Tr_{D/\QQ}(z^{2})>0$, and $\Tr_{D/\QQ}$ restricted to $F_{0}$ is
$[D:F_{0}]\cdot\Tr_{F_{0}/\QQ}$, so $\Tr_{F_{0}/\QQ}(z^{2})>0$ for every
nonzero $z\in F_{0}$. Suppose $F_{0}$ had a complex place $\sigma$. The
embedding of $F_{0}$ into the product of its archimedean completions has dense
image, so there is $z\in F_{0}$ with $\sigma(z)$ within $\tfrac1{10}$ of $i$
and $|\rho(z)|<\varepsilon$ for every other archimedean place $\rho$. Then
$\Tr_{F_{0}/\QQ}(z^{2})=\sum_{\rho}\rho(z)^{2}$, where a complex place
contributes twice the real part of $\rho(z)^{2}$; the term from $\sigma$ is at
most $2\cdot(-1+\tfrac{1}{2})<-1$ and the remaining terms are bounded by
$2[F_{0}:\QQ]\varepsilon^{2}$, so the total is negative for $\varepsilon$ small.
This contradicts positivity, so every place of $F_{0}$ is real.
\end{proof}

\begin{theorem}[Albert \cite{Alb34,Alb35a,Alb35b}]\label{thm:albert}
Let $(D,\dagger)$ be a division algebra over $\QQ$ with a positive involution,
$F$ its centre, $F_{0}$ the subfield fixed by $\dagger$, $e_{0}=[F_{0}:\QQ]$
and $[D:F]=d^{2}$. Then $F_{0}$ is totally real, $[F:F_{0}]\in\{1,2\}$, and
exactly one of the following holds.
\begin{enumerate}[label=\textup{Type \Roman*.},leftmargin=4.2em,itemsep=2pt]
\item $D=F=F_{0}$ is a totally real field and $\dagger$ is the identity.
\item $F=F_{0}$ is totally real and $D$ is a quaternion algebra over $F$ with
  $D\otimes_{F,\rho}\RR\cong M_{2}(\RR)$ at every real place $\rho$;
  $\dagger$ is of orthogonal type.
\item $F=F_{0}$ is totally real and $D$ is a quaternion algebra over $F$ with
  $D\otimes_{F,\rho}\RR\cong\mathbb{H}$ at every real place $\rho$; $\dagger$
  is the canonical involution.
\item $F$ is a CM field with maximal totally real subfield $F_{0}$, and
  $\dagger$ induces complex conjugation on $F$.
\end{enumerate}
If moreover $D=\End^{0}(A)$ for an abelian variety $A$ of dimension $g$ with
$\dagger$ the Rosati involution of a polarisation, then
\[
  \text{\rm I: } e_{0}\mid g,
  \qquad
  \text{\rm II: } 2e_{0}\mid g,
  \qquad
  \text{\rm III: } 2e_{0}\mid g,
  \qquad
  \text{\rm IV: } e_{0}d^{2}\mid g .
\]
\end{theorem}

We use \Cref{thm:albert} as a quoted classification; \Cref{lem:totallyreal}
proves the part of it that we appeal to directly, and the four types with the
divisibility conditions are in \cite[Ch.~9]{BL04}.

\begin{figure}[htbp]
\centering
  \includegraphics[width=0.92\textwidth]{fig_signature}
\caption{The signatures $(p,q)$ of an imaginary quadratic action, with
$p+q=\dim A$, drawn on the surface of height $pq$, which is the dimension of
the family of abelian varieties carrying that action, the unitary Shimura
variety of signature $(p,q)$. Each curve on the surface collects the
signatures available in one dimension, and on it the diagonal point is the
highest. The divisibility conditions of \Cref{thm:albert} put no
constraint on which of them occur when $\End^{0}(A)=K$, by
\Cref{prop:whichtypes}(i). By \Cref{prop:signatureforced} the Weil line
consists of Hodge classes exactly at the filled points on the ridge
$p=q$, and at every open point it is a rational sub-Hodge structure of type
$\{(p,2n-p),(2n-p,p)\}$ carrying no Hodge class at all. The diagonal is the
locus that \Cref{cor:weilalgebraic} shows to be an algebraic subvariety of the
moduli.}
\label{fig:albert}
\end{figure}

\subsection{The polarisation condition is automatic}

\begin{lemma}[Rosati on an imaginary quadratic subfield]\label{lem:rosatiauto}
Let $A$ be an abelian variety with a polarisation $\eta$ and Rosati involution
$\dagger$, and let $\iota\colon K\hookrightarrow\End^{0}(A)$ be an embedding of
an imaginary quadratic field such that $\iota(K)$ is stable under $\dagger$.
Then
\[
  \iota(\tau)^{\dagger}=\iota(\bbar{\tau})
  \qquad\text{for every }\tau\in K .
\]
That is, condition \textup{(iii)} of \Cref{def:weiltype} holds automatically,
and Weil type is the conjunction of \textup{(i)}, \textup{(ii)} and the
$\dagger$-stability of $\iota(K)$.
\end{lemma}

\begin{proof}
By stability, $\dagger$ restricts to a $\QQ$-algebra involution of the field
$\iota(K)\cong K$. A field automorphism of $K$ of order dividing two is either
the identity or complex conjugation, since $\Gal(K/\QQ)$ has order two and an
anti-automorphism of a commutative ring is an automorphism. Suppose it were
the identity, and put $x=\iota(\sqrt{-d})$. Then
$x\,x^{\dagger}=x^{2}=\iota(-d)=-d\cdot 1$, so
\[
  \Tr_{D/\QQ}\bigl(x\,x^{\dagger}\bigr)=-d\cdot\dim_{\QQ}D<0,
\]
contradicting positivity of $\dagger$. Hence $\dagger$ restricts to complex
conjugation.
\end{proof}

\begin{remark}\label{rem:whystable}
Stability is a genuine hypothesis and not a triviality: an arbitrary
subfield of $\End^{0}(A)$ need not be $\dagger$-stable. What
\Cref{lem:rosatiauto} removes is the second half of the condition, the
identification of $\dagger|_{K}$ with the specific automorphism. In
particular, in the model of \Cref{setup:model} one does not have to choose the
polarisation to match the field; any $\dagger$-stable embedding does the same
work.
\end{remark}

\subsection{The signature, and why balance is forced}

\begin{definition}[Signature]\label{def:signature}
Let $\iota\colon K\hookrightarrow\End^{0}(A)$ be as above, $g=\dim A$. The
operator $\iota(\sqrt{-d})$ acts on the tangent space $T_{0}A\cong\CC^{g}$ with
the two eigenvalues $\pm i\sqrt{d}$; write $p$ and $q=g-p$ for their
multiplicities. The pair $(p,q)$ is the \emph{signature} of $(A,\iota)$.
\end{definition}



The signature is an isogeny invariant, since an isogeny induces an isomorphism
on tangent spaces commuting with the action, and it is locally constant in
families, being the dimension of a kernel of a fixed operator on a bundle. By
duality $\iota(\tau)^{*}$ acts on $H^{1,0}(A)$ with the same multiplicities, so
$\dim V^{1,0}_{+}=p$ and $\dim V^{1,0}_{-}=q$ in the notation of
\Cref{sec:weiltype}.

\begin{proposition}[The balanced signature is forced]\label{prop:signatureforced}
Let $A$ be an abelian variety of dimension $2n$ with a $\dagger$-stable
embedding $\iota\colon K\hookrightarrow\End^{0}(A)$ of signature $(p,q)$,
$p+q=2n$, and let
\[
  \HW(A)=\Bigl(\textstyle\bigwedge^{2n}_{\CC}V_{+}
   \oplus\bigwedge^{2n}_{\CC}V_{-}\Bigr)\cap H^{2n}(A,\QQ)
\]
as in \Cref{def:weilline}. Then $\HW(A)$ is a rational sub-Hodge structure of
$H^{2n}(A,\QQ)$ of dimension two and of Hodge type
$\{(p,2n-p),\,(2n-p,p)\}$, and the following are equivalent:
\begin{enumerate}[label=\textup{(\roman*)},nosep]
\item $p=q=n$, that is $(A,\iota,\eta)$ is of $(K,-1,n)$-Weil type;
\item $\HW(A)$ contains a nonzero Hodge class;
\item $\HW(A)$ consists of Hodge classes.
\end{enumerate}
\end{proposition}

\begin{proof}
The space $\bigwedge^{2n}V_{+}$ is a line, and by \eqref{eq:bisplit} it equals
$\bigwedge^{p}V^{1,0}_{+}\otimes\bigwedge^{2n-p}V^{0,1}_{+}$, which is of pure
Hodge bidegree $(p,2n-p)$; the conjugate line $\bigwedge^{2n}V_{-}$ is of
bidegree $(2n-p,p)$. Their sum is stable under conjugation and defined over
$\QQ$ by the argument of \Cref{prop:weilline}(i), which never used the
balanced condition, so $\HW(A)$ is a two-dimensional rational sub-Hodge
structure of the stated type.

A Hodge class is an element of bidegree $(n,n)$. A nonzero element of
$\HW(A)\otimes\CC$ has a nonzero component in one of the two lines, so it is of
bidegree $(n,n)$ only if $(p,2n-p)=(n,n)$ or $(2n-p,p)=(n,n)$, both of which
say $p=n$. This gives (ii) implies (i). If $p=n$ then both lines are of
bidegree $(n,n)$ and every element of $\HW(A)$ is a Hodge class, which is (i)
implies (iii); and (iii) implies (ii) because $\dim\HW(A)=2>0$.
\end{proof}

\begin{remark}\label{rem:onlybalanced}
\Cref{prop:signatureforced} is the reason the literature speaks of abelian
varieties of Weil type rather than of imaginary quadratic multiplication in
general. For an unbalanced signature the space $\HW(A)$ still exists, is still
a canonical two-dimensional rational sub-Hodge structure, and is still not in
the divisor subring, but it contains no Hodge class and therefore predicts
nothing. The exceptional Hodge classes appear exactly at balance, and the
locus of balance is the algebraic subvariety of \Cref{cor:weilalgebraic}. The
statement is verified over all signatures $p$ and all $n\le3$ in
\texttt{explicit\_weil.py}, check (W4).
\end{remark}

\subsection{Which algebras can carry Weil classes}

\begin{proposition}[Weil type in the classification]\label{prop:whichtypes}
Let $A$ be an abelian variety of dimension $2n$ admitting a $\dagger$-stable
embedding of an imaginary quadratic field $K$ with balanced signature, and
suppose $\End^{0}(A)$ is a division algebra $D$ with centre $F$ as in
\Cref{thm:albert}. Then:
\begin{enumerate}[label=\textup{(\roman*)},nosep]
\item if $D=K$ then $D$ is of type IV with $e_{0}=1$ and $d=1$, and the
  divisibility condition of \Cref{thm:albert} is vacuous, so every even
  dimension occurs;
\item $D$ can also be of type III with $F_{0}=\QQ$, since a totally definite
  quaternion algebra over $\QQ$ contains imaginary quadratic subfields on which
  the canonical involution is conjugation;
\item $D$ cannot be of type I unless $D\subsetneq\End^{0}(A)$, because a
  totally real field contains no imaginary quadratic subfield.
\end{enumerate}
\end{proposition}

\begin{proof}
(i) If $D=K$ then $F=K$ is a CM field with $F_{0}=\QQ$, so $e_{0}=1$; and
$[D:F]=1$ gives $d=1$. The involution is conjugation by
\Cref{lem:rosatiauto}, which is type IV. The condition $e_{0}d^{2}\mid g$ reads
$1\mid 2n$.

(ii) Let $D$ be a quaternion algebra over $\QQ$ ramified at the infinite place,
with standard generators $u,v$ satisfying $u^{2}=a<0$, $v^{2}=b<0$, $uv=-vu$,
and canonical involution $x\mapsto\bbar{x}$ fixing $\QQ$ and negating
$u,v,uv$. Then $\QQ(u)\cong\QQ(\sqrt{a})$ is imaginary quadratic, stable under
the involution, and the involution restricts to conjugation on it. Whether
such a $D$ actually occurs as $\End^{0}(A)$ together with a balanced signature
is a separate question about the moduli, but nothing in the classification
excludes it.

(iii) Immediate, a totally real field has no element of negative square.
\end{proof}

\begin{remark}[The generic case is the one that matters]\label{rem:genericcase}
Case (ii) of \Cref{prop:whichtypes} is a special position: there
$\End^{0}(A)$ is larger than $K$, the Hodge group is correspondingly smaller,
and $\Hdg^{n}(A)$ can be larger than the three dimensions of
\Cref{thm:hdgcount}. It also tends to be the case where the Weil classes are
accessible; \Cref{rem:special} records that at such points the Weil line can
meet the divisor subring, which is the split situation. The hard case, and the
one all of \Cref{sec:divisor,sec:characters,sec:hyperplane} address, is
case (i) with $\Hg(A)=\SU(V,H)$.
\end{remark}

\begin{remark}[Why dimension four is where the subject starts]\label{rem:dimfour}
For $n=1$ the Weil line lies in $H^{2}(A,\QQ)$ and consists of divisor
classes, algebraic by the Lefschetz theorem on $(1,1)$ classes, so it predicts
nothing new; this is the same boundary as in \Cref{prop:normpart}. Moonen and
Zarhin determined the Hodge classes on abelian varieties of dimension at most
five \cite{MZ99}, and on simple abelian fourfolds in \cite{MZ95}: through
dimension three the Hodge ring is generated by divisor classes, and the first
Hodge classes not of that form occur in dimension four, on abelian fourfolds of
Weil type. So $n=2$ is simultaneously the first case of
\Cref{def:weiltype} with anything in it, the first case where
\Cref{thm:hdgcount} gives three rather than one, and the case that was the
first open one of the Hodge conjecture for abelian varieties until Markman
settled it \cite[Corollary 1.6.1]{Mar25a}; the first open case is now $n=3$
at nontrivial discriminant, by \Cref{prop:separate}(ii).
\end{remark}
